\subsection{Actual dominance derived from an original prescribed-base point}

\begin{lemma}[A faithful prescribed point forces actual dominance]
\label{berkovich:actual-prescribed-point-dominance}
\lean{EnvelopeProperty.Berkovich.RelativeSchemeValuedPoint.isDominant_of_prescribedPoint}
\leanok
If an actual scheme over an original domain has an actual relative native point
whose prescribed coefficient valuation has zero support ideal, its original
structure map is dominant. No dominance or generic-support certificate is supplied.
\end{lemma}
\begin{proof}
The genuine original affine-point equivalence identifies the base point's
actual prime with the original valuation support ideal, hence with the generic
prime. The actual relative point equality therefore puts the actual affine
base generic point in the original structure map's range. Its singleton has
closure the whole original affine spectrum. Thus the original range is dense,
which is precisely actual scheme dominance.
\end{proof}

\begin{lemma}[An actual positive-radius trivial Gauss point forces dominance]
\label{berkovich:actual-gauss-point-dominance}
\lean{EnvelopeProperty.Berkovich.RelativeSchemeValuedPoint.isDominant_of_gaussPoint}
\leanok
\uses{berkovich:actual-prescribed-point-dominance}
Over an original polynomial ring over any field, existence of an actual native
point over the original trivial-coefficient Gauss valuation of any positive
radius derives dominance of the SAME original structure morphism.
\end{lemma}
\begin{proof}
Every nonzero original polynomial has a nonzero original coefficient.
The corresponding weighted term is strictly positive at positive radius,
and is bounded above by the original Gauss value. Thus the actual Gauss
valuation has zero support. Apply the preceding whole native-point theorem.
\end{proof}
